% =====================================================================
% 11. Conclusiones
% =====================================================================
\section{Conclusiones}
\label{sec:conclusiones}

El proyecto entrega a la UPME una metodología reproducible, construida sobre
MongoDB, para contar las edificaciones y estimar el área de techos de los 170
municipios PDET con dos fuentes abiertas. El flujo va de los límites oficiales
del DANE y las huellas de Google y Microsoft hasta una colección
\campo{resultados} con 680 documentos verificados contra los conteos de las
colecciones, y de ahí a tablas, mapas, un ranking y una estimación de potencia,
todo con un solo comando y desde una base vacía.

Las decisiones de base de datos sostienen el resultado: los índices
\campo{2dsphere} hacen que la asignación de 4,5 millones de edificios use
\campo{IXSCAN} y no recorra las colecciones; la regla del centroide evita
contar dos veces un techo en el límite; la agregación corre dentro de la base;
los validadores \campo{\$jsonSchema} y el índice único sobre \campo{geom\_hash}
garantizan la integridad; y \campo{fuentes} y \campo{bitacora\_cargas} dejan
trazado de dónde viene cada dato y con qué parámetros se calculó. La asignación
coincidió al 100\,\% con una implementación independiente en GeoPandas.

En cuanto a los resultados, las dos fuentes coinciden en el total (222 y
221\,km\textsuperscript{2}) y en el orden de los municipios, pero no en el
Pacífico, donde Google tiene huecos de detección. El área se concentra: 15
municipios, encabezados por Santa Marta, Valledupar y Buenaventura, suman el
39\,\% del área conservadora, con una potencia instalable del orden de
6\,GWp. Esos municipios son los candidatos naturales para los pilotos; los del
Pacífico requieren verificación antes de decidir, y cualquier dimensionamiento
debe refinar el área útil y cruzarla con la irradiación y la demanda local.
